% !TeX root = part01.tex

\documentclass[../final.tex]{subfiles}

\begin{document}


\chapter{Взаимодействие частиц с веществом}

\rsubtitle{Ионизационные потери энергии}

% ============================================================
\section{Инвариантная масса}
% ============================================================

При анализе процессов взаимодействия частиц удобно использовать
инвариантную массу системы. Она определяется выражением

\begin{equation}
    M_{\mathrm{inv}}^2 =
    \left(\sum_i E_i\right)^2 -
    \left(\sum_i \vec{p}_i\right)^2 .
\end{equation}

Данная величина является лоренц-инвариантной, то есть её значение не
зависит от выбора системы отсчёта.

Инвариантная масса широко используется при поиске новых частиц и
резонансных состояний по распределению числа событий в зависимости от
массы системы.

% TODO: вставить рисунок из лекции:
% Распределение числа событий от M_inv.
% Пример реакции:
% p+p -> p+p+\pi^+\pi^-\pi^0


% ============================================================
\section{Взаимодействие частиц с веществом}
% ============================================================

При прохождении частицы через вещество происходит передача энергии
атомам среды. Основные механизмы взаимодействия:

\begin{itemize}
    \item ионизация атомов вещества;
    \item возбуждение электронных оболочек атомов;
    \item упругое рассеяние на ядрах;
    \item радиационные потери энергии.
\end{itemize}

Для заряженных частиц основным механизмом потери энергии при умеренных
энергиях является ионизация.


% ============================================================
\section{Ионизационные потери}
% ============================================================

Рассмотрим столкновение заряженной частицы с электроном вещества.
Максимальная энергия, которая может быть передана электрону отдачи,
определяется выражением

\begin{equation}
    E_{\max}
    =
    \frac{
        2m_e c^2 \beta^2 \gamma^2
    }
    {
        1+2\gamma m_e/M+(m_e/M)^2
    } ,
\end{equation}

где:

\begin{itemize}
    \item $M$ --- масса налетающей частицы;
    \item $m_e$ --- масса электрона;
    \item $\beta=v/c$;
    \item $\gamma=1/\sqrt{1-\beta^2}$.
\end{itemize}


В зависимости от энергии частицы возможны различные приближения.


Для нерелятивистского случая:

\[
\gamma \approx 1 ,
\]

и максимальная энергия электрона отдачи имеет порядок

\begin{equation}
    E_{\max}\sim \beta^2 .
\end{equation}


Для ультрарелятивистских частиц:

\[
\gamma\gg1 ,
\]

энергия электрона отдачи возрастает как

\begin{equation}
    E_{\max}\sim E .
\end{equation}


% ============================================================
\section{Энергетические потери при прохождении слоя вещества}
% ============================================================

Пусть частица проходит через слой вещества толщины $\Delta x$.
Количество энергии, потерянное частицей, можно оценить как

\begin{equation}
    \Delta W=
    \sigma \Delta x n \varepsilon ,
\end{equation}

где

\begin{itemize}
    \item $\sigma$ --- сечение взаимодействия;
    \item $n$ --- концентрация атомов вещества;
    \item $\varepsilon$ --- средняя энергия, передаваемая за одно
    столкновение.
\end{itemize}


Концентрация атомов выражается через плотность вещества:

\begin{equation}
    n=
    \frac{N_A\rho}{A},
\end{equation}

где $A$ --- атомная масса вещества, $\rho$ --- его плотность.


Тогда:

\begin{equation}
    \Delta W
    =
    \sigma\Delta x
    \frac{N_A\rho}{A}
    \varepsilon .
\end{equation}


Следовательно, удельные потери энергии:

\begin{equation}
    \frac{dE}{dx}
    =
    \sigma
    \frac{N_A\rho}{A}
    \varepsilon .
\end{equation}


% ============================================================
\section{Рассеяние на ядрах}
% ============================================================

При взаимодействии заряженной частицы с ядром возникает кулоновское
рассеяние.

Дифференциальное сечение рассеяния Резерфорда имеет вид:

\begin{equation}
    \frac{d\sigma}{d\Omega}
    =
    \frac{Z^2 e^4}
    {4E^2}
    \frac{1}
    {\sin^4(\theta/2)} .
\end{equation}


Здесь:

\begin{itemize}
    \item $Z$ --- заряд ядра;
    \item $E$ --- энергия налетающей частицы;
    \item $\theta$ --- угол рассеяния.
\end{itemize}


% TODO: вставить рисунок:
% схема рассеяния частицы на ядре,
% угол theta и траектория частицы

% ============================================================
\section{Формула Мотта}
% ============================================================

Формула Резерфорда не учитывает спин электрона и релятивистские эффекты.
При учёте этих поправок получается формула Мотта.

Дифференциальное сечение рассеяния принимает вид

\begin{equation}
    \frac{d\sigma}{d\Omega}
    =
    \frac{d\sigma_R}{d\Omega}
    \left(
        1-\beta^2\sin^2\frac{\theta}{2}
    \right),
\end{equation}

где

\[
    \beta=\frac{v}{c}
\]

характеризует релятивистский режим движения частицы.


Поправочный множитель

\[
    1-\beta^2\sin^2\frac{\theta}{2}
\]

описывает влияние спина частицы и становится существенным при больших
энергиях.


% TODO: вставить рисунок:
% Сравнение распределений Резерфорда и Мотта


% ============================================================
\section{$\delta$-электроны}
% ============================================================

При столкновении быстрой заряженной частицы с электроном атома последний
может получить энергию, достаточную для самостоятельной ионизации
вещества.

Такие электроны называются $\delta$-электронами.


Дифференциальное распределение числа выбитых электронов имеет вид

\begin{equation}
    \frac{d^2N}{d\varepsilon dx}
    =
    C
    \frac{Z}{A}
    \frac{1}{\beta^2}
    \frac{1}{\varepsilon^2}
    \left(
        1-\beta^2\frac{\varepsilon}{E_{\max}}
    \right),
\end{equation}

где

\begin{itemize}
    \item $\varepsilon$ --- энергия электрона отдачи;
    \item $E_{\max}$ --- максимальная энергия передачи;
    \item $C$ --- постоянная, зависящая от системы единиц.
\end{itemize}


Из выражения видно, что вероятность рождения электронов с малой энергией
значительно выше, так как

\[
    \frac{dN}{d\varepsilon}\sim\frac1{\varepsilon^2}.
\]


% ============================================================
\section{Средние ионизационные потери}
% ============================================================

Средняя потеря энергии частицы при прохождении вещества описывается
формулой Бете--Блоха.

Общий вид:

\begin{equation}
    -\frac{dE}{dx}
    =
    K
    \frac{Z}{A}
    \frac{z^2}{\beta^2}
    \left[
        \frac12
        \ln
        \frac{
            2m_ec^2\beta^2\gamma^2E_{\max}
        }
        {I^2}
        -
        \beta^2
        -
        \frac{\delta}{2}
    \right],
\end{equation}

где:


\begin{itemize}
    \item $z$ --- заряд налетающей частицы в единицах $e$;
    \item $Z$ --- заряд ядра вещества;
    \item $A$ --- атомная масса вещества;
    \item $I$ --- средний потенциал ионизации;
    \item $\delta$ --- поправка на поляризацию среды;
    \item $K$ --- константа Бете.
\end{itemize}


Константа:

\begin{equation}
    K=
    4\pi N_A r_e^2m_ec^2
    \approx
    0.307
    \frac{\mathrm{MeV\,cm^2}}{\mathrm{g}} .
\end{equation}


Формула Бете--Блоха показывает основные особенности потерь энергии:

\begin{itemize}
    \item при малых скоростях потери растут как $1/\beta^2$;
    \item при больших энергиях возникает логарифмический рост;
    \item существует область минимальных потерь.
\end{itemize}


% TODO: вставить график:
% Зависимость -dE/dx от импульса или энергии частицы


% ============================================================
\section{Минимально ионизирующие частицы}
% ============================================================

При определённой скорости частицы функция

\[
    -\frac{dE}{dx}
\]

имеет минимум.

Частицы в этой области называются минимально ионизирующими
частицами (MIP --- Minimum Ionizing Particle).


Для большинства материалов минимум достигается примерно при

\[
    \beta\gamma \approx 3-4 .
\]


В этой области потеря энергии составляет приблизительно

\[
    \frac{dE}{dx}
    \approx
    1-2
    \frac{\mathrm{MeV\,cm^2}}{\mathrm{g}} .
\]

% ============================================================
\section{Флуктуации потерь энергии}
% ============================================================

Среднее значение потери энергии, описываемое формулой Бете--Блоха,
не полностью характеризует процесс взаимодействия частицы с веществом.

Причина заключается в том, что число столкновений конечно, поэтому от
события к событию энергия, потерянная частицей, различается.


Распределение потерь энергии является асимметричным. Большинство событий
имеют небольшую потерю энергии, однако существует вероятность редких
столкновений с большой передачей энергии.


% TODO: вставить рисунок:
% Распределение потерь энергии (распределение Ландау)


Для тонких слоёв вещества распределение описывается функцией Ландау.


\section{Распределение Ландау}

При прохождении частицы через тонкий слой вещества распределение энергии
потерь имеет вид

\begin{equation}
    f(\Delta E)
    =
    \frac{1}{\xi}
    \phi(\lambda),
\end{equation}

где

\begin{equation}
    \lambda
    =
    \frac{
        \Delta E-\Delta E_{\mathrm{mp}}
    }
    {\xi}.
\end{equation}


Здесь:

\begin{itemize}
    \item $\Delta E_{\mathrm{mp}}$ --- наиболее вероятная потеря энергии;
    \item $\xi$ --- параметр, характеризующий ширину распределения;
    \item $\phi(\lambda)$ --- функция Ландау.
\end{itemize}


Параметр $\xi$ определяется как

\begin{equation}
    \xi
    =
    \frac{K}{2}
    \frac{Z}{A}
    \frac{x}{\beta^2},
\end{equation}

где $x$ --- толщина слоя вещества.


Наиболее вероятная потеря энергии:

\begin{equation}
    \Delta E_{\mathrm{mp}}
    =
    \xi
    \left[
        \ln
        \frac{2m_ec^2\beta^2\gamma^2}{I}
        +
        \ln
        \frac{\xi}{I}
        +
        j
        -
        \beta^2
        -
        \delta
    \right],
\end{equation}

где $j$ --- численная поправка.


% ============================================================
\section{Толстые слои вещества}
% ============================================================

При увеличении толщины вещества число столкновений возрастает и
распределение потерь энергии становится более симметричным.

В пределе большого числа независимых столкновений распределение
приближается к нормальному:

\begin{equation}
    f(E)
    =
    \frac{1}
    {\sqrt{2\pi\sigma^2}}
    \exp
    \left(
        -
        \frac{(E-\overline{E})^2}
        {2\sigma^2}
    \right).
\end{equation}


% TODO: вставить рисунок:
% Переход от распределения Ландау к гауссовскому


% ============================================================
\section{Радиационные потери энергии}
% ============================================================

Для лёгких заряженных частиц при больших энергиях существенную роль
начинают играть радиационные потери.

Основной процесс:

\[
    e^- \rightarrow e^-+\gamma .
\]

То есть частица теряет энергию за счёт тормозного излучения.


Интенсивность радиационных потерь пропорциональна энергии частицы:

\begin{equation}
    -\frac{dE}{dx}
    \sim E .
\end{equation}


В отличие от ионизационных потерь, которые зависят от скорости как
$1/\beta^2$, радиационные потери быстро возрастают с энергией.


Полная потеря энергии может быть записана в виде:

\begin{equation}
    -\frac{dE}{dx}
    =
    \left(
        \frac{dE}{dx}
    \right)_{\mathrm{ion}}
    +
    \left(
        \frac{dE}{dx}
    \right)_{\mathrm{rad}} .
\end{equation}


% ============================================================
\section{Радиационная длина}
% ============================================================

Для описания тормозного излучения вводится радиационная длина $X_0$.

После прохождения расстояния $X_0$ энергия электрона уменьшается в $e$
раз:

\begin{equation}
    E(x)=E_0e^{-x/X_0}.
\end{equation}


Радиационная длина зависит от материала и определяется его атомным
номером и плотностью.


% TODO: вставить таблицу:
% Радиационные длины различных материалов


% ============================================================
\section{Итог}
% ============================================================

Основные механизмы взаимодействия заряженных частиц с веществом:

\begin{enumerate}
    \item Ионизационные потери --- основной процесс при умеренных энергиях.
    \item Рассеяние на ядрах --- изменение направления движения частицы.
    \item Радиационные потери --- основной процесс для лёгких частиц при
    высоких энергиях.
\end{enumerate}


Для экспериментальной физики частиц наиболее важными являются:

\begin{itemize}
    \item формула Бете--Блоха для определения потерь энергии;
    \item распределение Ландау для описания флуктуаций;
    \item радиационная длина для описания электромагнитных каскадов.
\end{itemize}


\end{document}